\subsection{Exercises}

\begin{exercise}
    Let $X$ be an infinite set. Show the following statements without using \axiomsymbol{AC}:
    \begin{intheoremenumerate}
        \item Let $Y$ be a set with $X\subseteq Y$, then $Y$ is infinite.
        \item Let $Y$ be a set and assume there exists an injection $f\colon X\to Y$. Then $Y$ is infinite.
        \item $\powerset(X)$ is infinite.
    \end{intheoremenumerate}
\end{exercise}

\begin{exercise}
    Let $X$ be a Dedekind-infinite set. Show the following statements without using \axiomsymbol{AC}:
    \begin{intheoremenumerate}
        \item Let $Y$ be a set with $X\subseteq Y$, then $Y$ is Dedekind-infinite.
        \item For every finite subset $A\subseteq X$ holds that $X\setminus A$ is Dedekind-infinite.
        \item $\powerset(X)$ is Dedekind-infinite.
    \end{intheoremenumerate}
\end{exercise}

\begin{exercise}\label{exercise:finite_infinite:injection_from_infinite}
    Let $X$ and $Y$ be sets, $X$ infinite and $f\colon X\to Y$ an injection. Show that $Y$ is also infinite.
\end{exercise}
\begin{solution}
    By theorem \ref{theorem:finite_infinite:infinite_iff_injections} there exists an injection $g\colon\setN\to X$. Then $f\circ g\colon\setN\to Y$ is also an injection. Hence $Y$ is infinite by theorem \ref{theorem:finite_infinite:infinite_iff_injections}.
\end{solution}

\begin{exercise}
    Let $X$ be a set. Show that if $X$ is infinite then $\powerset(X)$ is infinite.
\end{exercise}
\begin{solution}
    The singleton mapping $s\colon X\to\powerset(X), x\mapsto\{x\}$ is an injection (see example \ref{example:mappings:singleton_mapping}). Therefore $\powerset(X)$ is infinite by exercise \ref{exercise:finite_infinite:injection_from_infinite}.
\end{solution}

\begin{exercise}\label{exercise:finite_infinite:pigeonhole_principle}
    Give a proof of the Pingeonhole Principle without the Bernstein-Cantor-Schröder theorem. Hint: Use the induction axiom and consider the cases $f(\setN_m)\subseteq\setN_n$ and $f(\setN_m)\nsubseteq\setN_n$ for an injection $f\colon\setN_m\to\setN_n$ to produce a contradiction.
\end{exercise}

\begin{solution}
    We proceed by induction on $n$:

    \proofstep{Base Case: $n=0$} For the base case $n=0$, let $m > 0$. Then $\setN_m \neq \emptyset$, whereas $\setN_0 = \emptyset$. Since there exists no mapping from a non-empty set to the empty set, there can be no injection $f \colon \setN_m \to \setN_0$.
    
    \proofstep{Inductive Step: $n \Rightarrow n+1$} Let $n \in \setN$ and assume the inductive hypothesis: for all $m > n$, there exists no injection from $\setN_m$ to $\setN_n$. Now, let $m \in \setN$ such that $m > n+1$. To reach a contradiction, assume there exists an injection $f \colon \setN_m \to \setN_{n+1}$.
    
    \proofstep{Case 1: $f(\setN_m) \subseteq \setN_n$}If the image of $f$ is contained within $\setN_n$, then the restricted mapping $f' \colon \setN_m \to \setN_n, b \mapsto f(b)$, is a well-defined injection. However, this directly contradicts the inductive hypothesis, as $m > n$.
    
    \proofstep{Case 2: $f(\setN_m) \not\subseteq \setN_n$} In this case, there must exist an element $a \in \setN_m$ such that $f(a) \notin \setN_n$. Since the codomain is $\setN_{n+1}$, it follows that $f(a) \in \setN_{n+1} \setminus \setN_n = \{n+1\}$, so $f(a) = n+1$. Furthermore, because $f$ is injective, $a$ is the unique element for which $f(a) = n+1$. Consequently, $f(b) \in \setN_n$ for all $b \in \setN_m \setminus \{a\}$.
    
    We now define the auxiliary mapping: \begin{align}g \colon \setN_{m-1} \rightarrow \setN_m, \qquad b \mapsto \begin{cases} b, & b < a \\ b+1, & b \geq a \end{cases}\end{align} Note that $g$ is injective and its image is $\setN_m \setminus \{a\}$. Thus, $f \circ g \colon \setN_{m-1} \to \setN_{n+1}$ is an injection, as the composition of two injective mappings is itself injective (see exercise \ref{exercise:mappings:composition_injections}). Since $a \notin g(\setN_{m-1})$ and $a$ is the only element in $\setN_n$ with $f(a)=n+1$, it follows that $(f \circ g)(\setN_{m-1}) \subseteq \setN_n$. Therefore, the restricted mapping \begin{align}h \colon \setN_{m-1} \rightarrow \setN_n,\qquad b \mapsto (f \circ g)(b),\end{align} is a well-defined injection. However, since $m > n+1$ implies $m-1 > n$, the existence of $h$ contradicts the inductive hypothesis.
    
    Having reached a contradiction in both cases, we conclude that no injection exists from $\setN_m$ to $\setN_{n+1}$ for $m > n+1$.
\end{solution}